% 英文源分章 PDF 第 18 页；第 9.10 节及致谢。
\section{讨论}

当实验和当前假设不足以将关注参数限定到可以采用传统的预先指定参数方法时，数据自适应参数这一一般概念，为估计相关科学参数带来了广阔机会。由于不预先指定关注参数会损失功效，因此必须谨慎地充分利用数据中所含的信息，估计自适应定义的参数，并构造有效的置信区间。CV-TMLE 实现了这两个目标。此外，该算法可以相对容易地并行化，基于影响曲线的推断也避免了耗时的自助法程序。因此，这一方法可以用于高维、大数据环境下的探索性数据分析。最后，在参数预先指定时，CV-TMLE 使估计量更接近完全自动化，因为此时就连用于估计数据生成分布的机器学习算法的自适应性问题，也不会影响其渐近性质，从而能够在最少假设下进行可靠推断。

由此，我们向这样一类统计算法又迈进了一步：它们只需要用户提供极少输入，就能在非常大的统计模型中产生相对有效的结果。再将数据自适应的参数定义纳入其中，自动化的潜力就更大了，因为连关注参数本身也不再需要事先（精确）指定。因此，这种面向数据自适应目标参数的 CV-TMLE 方法，是向经验不多的用户提供高性能探索性统计算法这一目标迈出的重要一步。我们正以这种方式走向一种情形：正如一个人即使不太了解机动车的机械原理，也仍然能够安全、有效地使用它。

\noindent\textbf{致谢。}本研究部分获得患者中心结局研究所（Patient-Centered Outcomes Research Institute，PCORI）试点项目计划资助（ME-1306-02735）。免责声明：本章全部陈述，包括研究发现和结论，仅代表作者本人，并不一定代表患者中心结局研究所（PCORI）、其理事会或方法学委员会的观点。
\clearpage
\endgroup
